% !PRINTED 375
\weberpage{375}
%<<CONT
Bestimmung von einer achten Doppeltangente aus einem gegebenen vollständigen Systeme durchzuführen. Wir bestimmen aber besser gleichzeitig drei. Es sei also jetzt

\weq{1}{x_1,\ x_2,\ x_3,\ x_4,\ x_5,\ x_6,\ x_7}

ein als bekannt vorausgesetztes vollständiges Siebener-System. Wir denken uns die drei \gesperrt{Steiner'schen} Complexe gebildet, in denen je sechs der Doppeltangenten (1), mit Ausschluss zuerst von $x_1$, dann von $x_2$, zuletzt von $x_3$ enthalten sind.

Die darin neu hinzutretenden 15 Doppeltangenten bezeichnen wir mit dem Buchstaben $\xi$ und erhalten diese drei Complexe [nach \S.\,99, (1), (2)] in der Gestalt:

\weq{2}{\begin{aligned}
x_2\xi_3 &,\ x_3\xi_2,\ x_4\xi_{41},\ x_5\xi_{51},\ x_6\xi_{61},\ x_7\xi_{71}\\
x_3\xi_1 &,\ x_1\xi_3,\ x_4\xi_{42},\ x_5\xi_{52},\ x_6\xi_{62},\ x_7\xi_{72}\\
x_1\xi_2 &,\ x_2\xi_1,\ x_4\xi_{43},\ x_5\xi_{53},\ x_6\xi_{63},\ x_7\xi_{73},
\end{aligned}}

und nach der Bezeichnung des \S.\,99 ist

\[
\xi_1 = [2\,3],\ \xi_2 = [3\,1],\ \xi_3 = [1\,2],\ \xi_{41} = [1\,4],\ \ldots
\]

Aus (2) ergiebt sich noch ein \gesperrt{Steiner'scher} Complex, der mit jedem der Complexe (2) ein syzygetisches Paar bildet:

\weq{3}{x_1\xi_1,\ x_2\xi_2,\ x_3\xi_3,\ \ldots}

und dieser Complex enthält keine der Doppeltangenten $x_4, x_5, x_6, x_7$.

Indem wir nun die gesuchten Functionen $\xi_1, \xi_2, \xi_3$ mit den geeigneten constanten Factoren multiplicirt annehmen, können wir die Gleichung der Curve vierter Ordnung nach \S.\,96, (9) in die Form

\weq{4}{\sqrt{x_1\xi_1} + \sqrt{x_2\xi_2} + \sqrt{x_3\xi_3} = 0}

setzen, und also die rationale Function $f$, die, gleich Null gesetzt, die Gleichung der Curve giebt, in jeder der drei mit einander identischen Formen

\weq{5}{f = 4x_2\xi_2x_3\xi_3 - u_1^2 = 4x_3\xi_3x_1\xi_1 - u_2^2 = 4x_1\xi_1x_2\xi_2 - u_3^2,}

annehmen, worin

\weq{6}{\begin{aligned}
u_1 &= -x_1\xi_1 + x_2\xi_2 + x_3\xi_3\\
u_2 &= x_1\xi_1 - x_2\xi_2 + x_3\xi_3\\
u_3 &= x_1\xi_1 + x_2\xi_2 - x_3\xi_3
\end{aligned}}

gesetzt ist. Nun bilden [nach (2)] auch $x_2\xi_3\, x_4\xi_{41}$ ein syzygetisches Quadrupel, und folglich können wir, wenn über einen
%CONT>>
